\documentclass{article}%
\usepackage[T1]{fontenc}%
\usepackage[utf8]{inputenc}%
\usepackage{lmodern}%
\usepackage{textcomp}%
\usepackage{lastpage}%
\usepackage{geometry}%
\geometry{margin=2.2cm}%
%
\usepackage{amsmath}%
\usepackage{amssymb}%
\usepackage{booktabs}%
\usepackage{longtable}%
\usepackage{array}%
\usepackage{hyperref}%
\hypersetup{hidelinks}%
\title{ITA 2026 {-} Química (claude\_opus\_5\_5) x Poliedro\\[4pt]\large Comparação com o gabarito Poliedro Resolve\thanks{Gerado com gpt-resolve: \url{https://github.com/lgabs/gpt-resolve}}}%
\date{10/11/2026}%
%
\begin{document}%
\normalsize%
\maketitle%
\section*{Resumo}\textbf{Nota média: 10,0/10} (33 de 33 itens com a resolução correta; cada item vale o mesmo dentro da questão).\par\medskip%
\begin{longtable}{@{}c c p{0.78\textwidth}@{}}\toprule Questão & Nota & Resumo \\ \midrule
1 & 10,0 & A resolução está correta e bate com o gabarito do Poliedro nos três itens: a) 0,4 mol/L; b) 2,0 mol/L; c) cerca de 2,1$\times$10$^{-}$$^2$ mol/L, já que 2/95 = 0,8/38. Os dois tratam a hidrólise de A$^{-}$ de forma exata, com Kb = 10$^{-}$$^2$ e metade do A$^{-}$ hidrolisada. Refiz as contas em python e o balanço de cargas fecha. A resolução também confere a consistência pelo balanço de cargas, e a resposta está completa. \\
2 & 10,0 & A resolução está correta nos dois itens e equivale ao gabarito do Poliedro. a) 0,612 L contra 0,6109 L. A diferença vem só de R: a resolução usa 0,0821, o valor dado na prova, e o Poliedro usa 0,082. b) 0,354 mol/L contra 0,353. A diferença vem de arredondamento: o valor exato é 0,3536, e o Poliedro calculou 0,707$^3$ em vez de 0,5$\cdot$0,707. O gabarito tem um deslize de unidade ('mol' em vez de mol/L), sem efeito no conteúdo. Nenhuma divergência real. \\
3 & 10,0 & A resolução está correta nos dois itens e concorda com o gabarito Poliedro, com diferenças só de arredondamento. No item a, o tempo exato é 552,8 s: a resolução dá 552 s (e cita $\approx$ 553 s), e o gabarito dá 553,6 s. Os dois variam só por onde aplicaram $\surd$5 $\approx$ 2,24. No item b, a resolução dá os valores exatos, {[}A+{]} $\approx$ 6,11$\cdot$10$^{-}$$^3$ e {[}B-{]} $\approx$ 1,64$\cdot$10$^{-}$$^3$ mol/L, cujo produto é exatamente o Kps. O gabarito chega a 6,135 e 1,655$\cdot$10$^{-}$$^3$ porque arredondou no meio da conta, e o produto desses valores fica 1,5\% acima do Kps. Não há erro de raciocínio em nenhum dos dois, e a resolução é numericamente mais precisa. \\
4 & 10,0 & A resolução está correta nos dois itens. No item b, os dois dão n = 9, e a conta da resolução confere: M\_sal = 466,35 g/mol, parte anidra = 303,00 g/mol, n = 9,06. No item a, a resolução dá 5,58 g e o gabarito dá 5,57 g, o que é equivalente. A resolução usou direto o resíduo medido de 28,5 \% (5,579 g). O gabarito usou a fórmula com n = 9 e massas molares arredondadas (5,565 g); a resolução também mostra esse valor (5,56 g). O gabarito tem um erro de digitação na equação do item b: falta o fator 2 no denominador, e a equação escrita levaria a n $\approx$ 35. A resposta final do gabarito, porém, está certa. \\
5 & 10,0 & A resolução da questão 5 está correta e bate com o gabarito do Poliedro Resolve em todos os itens. A identificação é igual (H, B, C, N, O, F, S). a) 3 arranjos para {[}CNS{]}$^{-}$. b) BF$_{3}$ (ou SO$_{3}$) é trigonal plana e NH$_{3}$ (ou NF$_{3}$) é piramidal, exemplos que estão na lista do gabarito. c) Nitração do fluorobenzeno com os produtos orto e para desenhados. A única diferença é que a resolução aponta o para como principal e o orto como minoritário, o que é quimicamente mais preciso, enquanto o gabarito chama os dois de 'principais'. d) H$_{2}$S < CH$_{3}$OH < H$_{2}$O. Não encontrei erro de conteúdo. As dúvidas listadas são pequenas: redação, leitura do enunciado e arredondamento. \\
6 & 10,0 & A resolução está correta e coincide com o gabarito do Poliedro nos três itens. Os pKa/pKb derivados também são iguais: pKb(CH3CC-) = -11, pKa(H3O+) = -1,70, pKa(NH4+) = 9,25 e pKa(NH3) = 36. As conclusões são: a) não favorece os produtos (K = 10\^{}-26,7); b) não favorece (K = 10\^{}-15,75); c) favorece (K = 10\^{}11). As contas foram conferidas numericamente. Os únicos pontos menores são a convenção do pKb da água, que os dois textos apontam, e uma aproximação 10\^{}1,70 $\approx$ 55,5 na observação da resolução, sem efeito na resposta. \\
7 & 10,0 & A resolução está correta e bate com o gabarito do Poliedro nos cinco itens. As únicas diferenças são numéricas e pequenas ($\leq$ 0,7\%). Vêm do gabarito ter arredondado N\_A para 6$\times$10$^2$$^3$ e M(H$_{2}$) para 2 g/mol, enquanto a resolução usa as constantes da prova: b) 532 contra 530,4 kJ; c) 1,35 contra 1,36$\times$10$^{-}$$^2$ mol; d) 1,49 contra 1,5 GJ; e) 837 contra 835 nm. Refiz todas as contas em Python e os valores da resolução estão corretos. No item a, as semirreações são equivalentes. A resolução usa o mesmo meio no par e dá as versões ácida e básica. O gabarito mistura meios, mas a global continua correta. Veredito: 5/5 itens equivalentes. \\
8 & 10,0 & A resolução da questão 8 está correta nos quatro itens e bate com o gabarito do Poliedro. Os itens a (-800 kJ/mol; 15 mol para -12 MJ) e d (200 g de CH4(l)) são idênticos. Nos itens b e c a diferença é só de arredondamento: a resolução usa M(CH4) = 16,05 g/mol, valor da tabela da prova, e obtém 240,75 g e 0,644 g; o gabarito usa 16 g/mol e obtém 240 g e 0,64 g. Refiz todas as contas em Python e elas conferem. O enunciado transcrito confere com a página 10 da prova. Nenhum ponto altera a nota. \\
9 & 10,0 & A resolução concorda com o gabarito Poliedro nos quatro itens. Em (a), a lista de espécies é a mesma; a resolução ainda acrescenta o HCrO$_{4}$$^{-}$ e classifica melhor o HSO$_{4}$$^{-}$, que é majoritário em H$_{2}$SO$_{4}$ diluído e que o gabarito põe entre os traços. Em (b), as duas explicam por migração iônica: o Cu$^2$? vai para Pt$_{1}$ e o Cr$_{2}$O$_{7}$$^2$$^{-}$ para Pt$_{2}$. Em (c), Pt$_{1}$ é o catodo (2H? + 2e$^{-}$ $\to$ H$_{2}$) e Pt$_{2}$ o anodo (2H$_{2}$O $\to$ O$_{2}$ + 4H? + 4e$^{-}$). Em (d), P$_{1}$ é negativo, P$_{2}$ é positivo e a corrente convencional vai de P$_{2}$ a P$_{1}$ pelo circuito externo. Não encontrei erro na resolução; o único ponto discutível é a classificação do HSO$_{4}$$^{-}$ no gabarito, que não muda a nota. \\
10 & 10,0 & A resolução está correta e bate com o gabarito do Poliedro nos três itens: a) X = 226, b) Z = 86, c) 3 partículas $\alpha$. O enunciado da página 12 do PDF original confere com a transcrição. O caminho é o mesmo nos dois: $\beta$$^{-}$ de Q{-}232 (Z = 89) dá G{-}232 (Z = 90, N = 142); E é isótono de G com Z = 84, logo tem massa 226; A é isóbaro de E; as contas de massa e carga dão n = 3 e Z = 86. A resolução ainda confere que A e E têm números atômicos diferentes. Não há divergência. \\
\bottomrule\end{longtable}%
\section*{Principais pontos de dúvida}%
\begin{itemize}%
\item%
\textbf{Q4} \textit{(key)} Erro de digitação no gabarito, item b. A equação está escrita como 0,285 = 266/(303 + 18n). Lida ao pé da letra, ela dá 303 + 18n = 933,3, logo n $\approx$ 35, e não 9. O próprio texto do gabarito diz que 2 mol de oxalato formam 1 mol de Nb2O5, então a forma correta é 0,285 = 266/{[}2(303 + 18n){]}, ou 133/(303 + 18n). A resposta final n = 9 está certa; o erro está só na equação escrita.%
\end{itemize}%
\section*{Questão 1}%
\href{https://poliedroresolve.sistemapoliedro.com.br/vestibulares/ita/2026/ita-2-fase-dia-2-29-10-2025-2-fase-dia-2/dissertativa-1-quimica-2-fase-dia-2-ita-2026}{Resolução do gabarito}\par\smallskip%
\begin{longtable}{@{}p{0.06\textwidth} p{0.33\textwidth} p{0.33\textwidth} p{0.1\textwidth} p{0.08\textwidth}@{}}\toprule Item & Resolução & Gabarito & Veredito & Certo \\ \midrule
a & 0,4 mol/L & 0,4 mol/L & igual & ambos \\
b & 2 mol/L & 2,0 mol/L & igual & ambos \\
c & 2/95 $\approx$ 2,1$\times$10$^{-}$$^2$ mol/L & $\approx$ 2,1$\cdot$10$^{-}$$^2$ mol/L (0,8/38) & equivalente & ambos \\
\bottomrule\end{longtable}%
\begin{itemize}%
\item%
\textbf{a)} Kb = Kw/Ka = 10$^{-}$$^2$. Na mistura final (40 mL), {[}OH$^{-}${]} = {[}HA{]} = 10$^{-}$$^2$, então Kb = 10$^{-}$$^{4}$/(C - 10$^{-}$$^2$) = 10$^{-}$$^2$, o que dá C = 2$\times$10$^{-}$$^2$ mol/L e n = 8$\times$10$^{-}$$^{4}$ mol, contidos em 2 mL. Logo C$_{2}$ = 0,4 mol/L. O gabarito chega ao mesmo valor pela equação 5$\cdot$(M\_A)$_{2}$ - 1 = 1. Refiz a conta em python e confirmei.%
\item%
\textbf{b)} Diluição de 4 mL para 20 mL: C$_{1}$ = 0,4$\cdot$20/4 = 2,0 mol/L. É o mesmo valor nos dois.%
\item%
\textbf{c)} 2/95 e 0,8/38 são a mesma fração, $\approx$ 0,02105 mol/L. A proporção é 1:1 no ponto de equivalência: n\_base = 8$\times$10$^{-}$$^{4}$ mol em 38 mL. O balanço de cargas fecha: {[}M?{]} = 2$\times$10$^{-}$$^2$ = {[}A$^{-}${]} + {[}OH$^{-}${]} = 10$^{-}$$^2$ + 10$^{-}$$^2$.%
\end{itemize}%
\paragraph{Pontos de dúvida}%
\begin{itemize}%
\item%
\textit{(statement)} O enunciado escreve "Solução I" na definição da Solução 2. É um erro de digitação, e o sentido é Solução 1. Tanto a resolução quanto o gabarito leram assim.%
\item%
\textit{(statement)} Com Ka = 10$^{-}$$^{1}$$^2$, a base conjugada tem Kb = 10$^{-}$$^2$, e metade do A$^{-}$ está hidrolisada no ponto de equivalência. Quem usasse a aproximação usual C $\approx$ {[}A$^{-}${]} chegaria a valores com metade do tamanho: 0,2, 1,0 e cerca de 1,05$\times$10$^{-}$$^2$ mol/L. A resolução e o gabarito fazem o cálculo exato, que é o correto. Vale observar em aula que essa aproximação não serve aqui.%
\end{itemize}%
\section*{Questão 2}%
\href{https://poliedroresolve.sistemapoliedro.com.br/vestibulares/ita/2026/ita-2-fase-dia-2-29-10-2025-2-fase-dia-2/dissertativa-2-quimica-2-fase-dia-2-ita-2026}{Resolução do gabarito}\par\smallskip%
\begin{longtable}{@{}p{0.06\textwidth} p{0.33\textwidth} p{0.33\textwidth} p{0.1\textwidth} p{0.08\textwidth}@{}}\toprule Item & Resolução & Gabarito & Veredito & Certo \\ \midrule
a & V(O2) = 0,025$\cdot$0,0821$\cdot$298 $\approx$ 0,612 L ($\approx$ 612 mL) & V(O2) $\approx$ 0,6109 L (0,025$\cdot$0,082$\cdot$298) & equivalente & ambos \\
b & {[}H2O2{]} = 0,5$\cdot$e\^{}(-0,3465) = 0,5$\cdot$0,707 $\approx$ 0,354 mol/L & {[}H2O2{]} = (e\^{}(-0,3465))\^{}3 = 0,707\^{}3 $\approx$ 0,353 mol/L (o gabarito escreveu 'mol') & equivalente & ambos \\
\bottomrule\end{longtable}%
\begin{itemize}%
\item%
\textbf{a)} O raciocínio é o mesmo: n(H2O2) = 0,5$\cdot$0,1 = 0,050 mol, logo n(O2) = 0,025 mol, e PV = nRT a 298 K e 1 atm. A diferença vem só do valor de R. A resolução usa R = 0,0821 atm$\cdot$L$\cdot$K$^{-}$$^{1}$$\cdot$mol$^{-}$$^{1}$, que é o valor dado na prova, e obtém 0,6116 $\approx$ 0,612 L. O gabarito arredonda R para 0,082 e obtém 0,6109 L. A diferença é de cerca de 0,1\%, ou seja, arredondamento. A resolução segue mais à risca a constante fornecida.%
\item%
\textbf{b)} Os dois usam cinética de 1ª ordem, indicada pela unidade min$^{-}$$^{1}$, com kt = 0,0693$\cdot$5 = 0,3465. A resolução calcula 0,5$\cdot$0,707 = 0,3535 $\approx$ 0,354. O gabarito escreve 0,5 = e\^{}(-0,693) = (e\^{}(-0,3465))\^{}2 e chega a 0,707\^{}3 = 0,35339 $\approx$ 0,353. O valor exato é 0,5$\cdot$e\^{}(-0,3465) = 0,35358, portanto 0,354 é o arredondamento mais fiel. A diferença na 3ª casa decimal vem só de arredondamento (0,707 $\approx$ 1/$\surd$2). A unidade 'mol' do gabarito é deslize de digitação: o pedido é uma concentração, em mol/L.%
\end{itemize}%
\paragraph{Pontos de dúvida}%
\begin{itemize}%
\item%
\textit{(key)} No item b, o gabarito dá a resposta em 'mol', mas o enunciado pede concentração, que deve estar em mol$\cdot$L$^{-}$$^{1}$. A passagem intermediária '-ln2 - 0,0693$\cdot$5 = -0,693 = -10$\cdot$0,0693' também está mal escrita, embora o resultado final esteja certo.%
\item%
\textit{(key)} No item a, o gabarito usa R = 0,082, mas a prova fornece R = 8,21$\times$10$^{-}$$^2$ atm$\cdot$L$\cdot$K$^{-}$$^{1}$$\cdot$mol$^{-}$$^{1}$. Por isso dá 0,6109 L em vez de 0,612 L. É só arredondamento.%
\item%
\textit{(statement)} O enunciado não diz se k se refere ao consumo de H2O2 (-d{[}H2O2{]}/dt = k{[}H2O2{]}) ou à velocidade da reação com o coeficiente 2 (-$\frac12$ d{[}H2O2{]}/dt = k{[}H2O2{]}). No segundo caso, o expoente efetivo seria 2kt e a resposta seria 0,25 mol/L. O dado e\^{}(-0,3465) = 0,707 da folha de constantes deixa claro que a intenção da banca é kt = 0,3465. Resolução e gabarito adotam essa leitura.%
\item%
\textit{(solution)} Hipóteses implícitas no item a: O2 tratado como gás ideal e seco, desprezando a solubilidade em água e a pressão de vapor da água. A resolução declara isso. São as hipóteses usuais e não mudam a resposta esperada.%
\end{itemize}%
\section*{Questão 3}%
\href{https://poliedroresolve.sistemapoliedro.com.br/vestibulares/ita/2026/ita-2-fase-dia-2-29-10-2025-2-fase-dia-2/dissertativa-3-quimica-2-fase-dia-2-ita-2026}{Resolução do gabarito}\par\smallskip%
\begin{longtable}{@{}p{0.06\textwidth} p{0.33\textwidth} p{0.33\textwidth} p{0.1\textwidth} p{0.08\textwidth}@{}}\toprule Item & Resolução & Gabarito & Veredito & Certo \\ \midrule
a & t $\approx$ 552 s (V* = $\surd$0,2 = $\surd$5/5 $\approx$ 0,448 L; evaporam 552 mL a 1 g/s; com $\surd$5 exato, $\approx$ 552,8 s $\approx$ 553 s) & $\Delta$t = 1240/2,24 s $\approx$ 553,6 s (V' = 1/2,24 L) & equivalente & ambos \\
b & {[}A+{]} = (1+$\surd$3)$\surd$5$\cdot$10$^{-}$$^3$ $\approx$ 6,1$\cdot$10$^{-}$$^3$ mol/L; {[}B-{]} = ($\surd$3-1)$\surd$5$\cdot$10$^{-}$$^3$ $\approx$ 1,6$\cdot$10$^{-}$$^3$ mol/L (V = V*/2 $\approx$ 0,224 L) & {[}A+{]} = 6,135$\cdot$10$^{-}$$^3$ mol/L; {[}B-{]} = 1,655$\cdot$10$^{-}$$^3$ mol/L & equivalente & ambos \\
\bottomrule\end{longtable}%
\begin{itemize}%
\item%
\textbf{a)} O raciocínio é o mesmo nos dois: a água evapora a 540/540 = 1 g/s = 1 mL/s, e a precipitação começa quando (2$\cdot$10$^{-}$$^3$/V)(1$\cdot$10$^{-}$$^3$/V) = 10$^{-}$$^{5}$, ou seja, V* = 1/$\surd$5 = 0,44721 L. Refazendo a conta, o tempo exato é t = 1000(1 - 1/$\surd$5) = 552,79 s. A resolução escreve V* = $\surd$5/5 e usa $\surd$5 = 2,24, o que dá 0,448 L e 552 s. O gabarito escreve V* = 1/2,24, o que dá 0,4464 L e 553,6 s. A diferença vem só de onde cada um aplicou o valor aproximado de $\surd$5: o exato fica entre os dois e a diferença é menor que 0,3\%. Além disso, a resolução informa também o valor 552,8 s $\approx$ 553 s.%
\item%
\textbf{b)} Os dois usam o mesmo modelo: V = V*/2, a precipitação ocorre na proporção 1:1 e no equilíbrio {[}A+{]}{[}B-{]} = Kps. Resolvendo exatamente (x$^2$ - 3x + 1,5 = 0, com x em 10$^{-}$$^3$ mol), a solução é {[}A+{]} = (1+$\surd$3)$\surd$5$\cdot$10$^{-}$$^3$ = 6,109$\cdot$10$^{-}$$^3$ e {[}B-{]} = ($\surd$3-1)$\surd$5$\cdot$10$^{-}$$^3$ = 1,637$\cdot$10$^{-}$$^3$ mol/L, e o produto dá exatamente 1,000$\cdot$10$^{-}$$^{5}$. Esse é o resultado da resolução. O gabarito chega a 6,135 e 1,655$\cdot$10$^{-}$$^3$ porque arredondou duas vezes no meio da conta: usou 8,96 e 4,48$\cdot$10$^{-}$$^3$ no lugar de 8,944 e 4,472$\cdot$10$^{-}$$^3$ e trocou o termo constante 3,014$\cdot$10$^{-}$$^{5}$ por 3$\cdot$10$^{-}$$^{5}$. Com isso, o produto das concentrações finais do gabarito dá 1,015$\cdot$10$^{-}$$^{5}$, cerca de 1,5\% acima do Kps, e {[}B-{]} fica cerca de 1,1\% acima do valor exato. O método está certo nos dois e a diferença é só de arredondamento. A resolução é mais exata, e com dois algarismos significativos {[}B-{]} = 1,6$\cdot$10$^{-}$$^3$, não 1,7$\cdot$10$^{-}$$^3$.%
\end{itemize}%
\paragraph{Pontos de dúvida}%
\begin{itemize}%
\item%
\textit{(key)} No item b, o gabarito arredonda o termo constante da equação de 2º grau (3,014$\cdot$10$^{-}$$^{5}$ para 3$\cdot$10$^{-}$$^{5}$) e parte de 8,96/4,48$\cdot$10$^{-}$$^3$ em vez de 8,944/4,472$\cdot$10$^{-}$$^3$. Por isso, as concentrações finais (6,135 e 1,655$\cdot$10$^{-}$$^3$) não satisfazem exatamente o Kps: o produto dá 1,015$\cdot$10$^{-}$$^{5}$. Os valores exatos são 6,109$\cdot$10$^{-}$$^3$ e 1,637$\cdot$10$^{-}$$^3$ mol/L. Com dois algarismos significativos, o {[}B-{]} do gabarito arredondaria para 1,7$\cdot$10$^{-}$$^3$, e o exato é 1,6$\cdot$10$^{-}$$^3$. Mesmo resolvendo a equação arredondada, sai x = 2,8266$\cdot$10$^{-}$$^3$ e não 2,825$\cdot$10$^{-}$$^3$ como o gabarito escreve.%
\item%
\textit{(solution)} No item a, a resposta principal da resolução (552 s) usa $\surd$5 = 2,24 dentro de $\surd$5/5, e o gabarito (553,6 s) usa 1/2,24. O valor exato é 552,8 s. A diferença é só de arredondamento: a resolução informa também o valor de \textasciitilde{}553 s, e o examinador deve aceitar qualquer valor próximo de 553 s.%
\item%
\textit{(statement)} A consideração ii (ebulição nas condições padrão, 1 bar) junto com a solução a 100 $^\circ$C é apenas nominal (a 1 bar a água ferve a \textasciitilde{}99,6 $^\circ$C). Os dois tratam todo o calor como vaporização a 100 $^\circ$C, com $\Delta$H = 540 cal/g. Também fica implícito que o volume da solução é igual ao volume de água líquida, inclusive depois que o AB(s) se forma, e os dois assumem isso.%
\end{itemize}%
\section*{Questão 4}%
\href{https://poliedroresolve.sistemapoliedro.com.br/vestibulares/ita/2026/ita-2-fase-dia-2-29-10-2025-2-fase-dia-2/dissertativa-4-quimica-2-fase-dia-2-ita-2026}{Resolução do gabarito}\par\smallskip%
\begin{longtable}{@{}p{0.06\textwidth} p{0.33\textwidth} p{0.33\textwidth} p{0.1\textwidth} p{0.08\textwidth}@{}}\toprule Item & Resolução & Gabarito & Veredito & Certo \\ \midrule
a & m $\approx$ 5,58 g de oxalato (1,590 g de Nb2O5 em 10 g de BGN, dividido por 0,285); a verificação com n = 9 dá 5,56 g & m $\approx$ 5,57 g de oxalato (usando n = 9, M = 465 g/mol, M média do BGN = 66,8 g/mol, n(Nb2O5) = 1/167 mol) & equivalente & ambos \\
b & n = 9 (M\_sal = 132,91/0,285 = 466,35 g/mol; M\_anidro = 303,00 g/mol; n = 163,35/18,02 = 9,06 $\approx$ 9) & n = 9 (2 mol de oxalato $\to$ 1 mol de Nb2O5; 0,285 = 266/(303 + 18n)) & igual & ambos \\
\bottomrule\end{longtable}%
\begin{itemize}%
\item%
\textbf{a)} Refiz as contas. Em 100 mol de BGN há 6687,56 g, dos quais 1063,28 g são Nb2O5, ou seja, w = 15,90 \%. Em 10 g de BGN há 1,5899 g de Nb2O5. Usando direto o resíduo medido: 1,5899/0,285 = 5,579 g. Esse é o valor da resolução. Usando a fórmula estequiométrica com n = 9 (M = 465,18 g/mol): 1,5899/265,82 $\cdot$ 2 $\cdot$ 465,18 = 5,565 g. O gabarito, com massas molares arredondadas, obtém 2$\cdot$465/167 = 5,57 g. A diferença vem só de usar o dado experimental de 28,5 \% ou a fórmula com n inteiro, que prevê 28,57 \% de resíduo. Isso dá uma discrepância de 0,2 \%, menor que a precisão do dado, e a própria resolução mostra os dois valores (5,58 e 5,56 g). As duas respostas são aceitáveis.%
\item%
\textbf{b)} Conferi as contas. A parte anidra NH4{[}NbO(C2O4)2{]} tem M = 14,01 + 4,04 + 92,91 + 16,00 + 2$\cdot$88,02 = 303,00 g/mol. A massa molar do sal sai de 265,82/2/0,285 = 466,35 g/mol, o que dá n = (466,35 - 303,00)/18,02 = 9,06 $\approx$ 9. Com n = 9 o resíduo teórico é 28,57 \%, contra 28,6 \% para n = 8 (29,7 \%) e n = 10 (27,5 \%) --- ou seja, só n = 9 bate com 28,5 \%. O valor final coincide com o do gabarito. A equação do gabarito, porém, está escrita sem o fator 2 (ver dúvidas).%
\end{itemize}%
\paragraph{Pontos de dúvida}%
\begin{itemize}%
\item%
\textit{(key)} Erro de digitação no gabarito, item b. A equação está escrita como 0,285 = 266/(303 + 18n). Lida ao pé da letra, ela dá 303 + 18n = 933,3, logo n $\approx$ 35, e não 9. O próprio texto do gabarito diz que 2 mol de oxalato formam 1 mol de Nb2O5, então a forma correta é 0,285 = 266/{[}2(303 + 18n){]}, ou 133/(303 + 18n). A resposta final n = 9 está certa; o erro está só na equação escrita.%
\item%
\textit{(key)} Item a: a resposta muda um pouco conforme o caminho. Usando direto o resíduo medido de 28,5 \% sai 5,58 g (resolução). Usando a fórmula com n = 9 (resíduo teórico de 28,57 \%) sai 5,56--5,57 g (gabarito). Os dois caminhos são válidos e a diferença (cerca de 0,2 \%) está dentro da precisão do dado. Pode valer confirmar com o professor qual valor a banca espera.%
\item%
\textit{(solution)} Hipóteses implícitas razoáveis, aceitas também pelo gabarito: a composição 60{-}36{-}4 é em mol (o enunciado diz 'porcentagem molar'); todo o Nb do precursor vai para o BGN; o único resíduo da termogravimetria é Nb2O5. O gabarito lista CO entre os gases e a resolução não, o que não afeta as contas.%
\end{itemize}%
\section*{Questão 5}%
\href{https://poliedroresolve.sistemapoliedro.com.br/vestibulares/ita/2026/ita-2-fase-dia-2-29-10-2025-2-fase-dia-2/dissertativa-5-quimica-2-fase-dia-2-ita-2026}{Resolução do gabarito}\par\smallskip%
\begin{longtable}{@{}p{0.06\textwidth} p{0.33\textwidth} p{0.33\textwidth} p{0.1\textwidth} p{0.08\textwidth}@{}}\toprule Item & Resolução & Gabarito & Veredito & Certo \\ \midrule
Identificação dos átomos & A = H, Z = B, X = C, G = N, R = O, Y = F, Q = S & A = H, Z = B, X = C, G = N, R = O, Y = F, Q = S & igual & ambos \\
a & 3 arranjos para {[}CNS{]}$^{-}$ (átomo central C, N ou S: S--C--N, C--N--S, C--S--N); o mais estável é o tiocianato {[}S--C$\equiv$N{]}$^{-}$ & 3 arranjos, conforme o átomo central: {[}CNS{]}$^{-}$, {[}NCS{]}$^{-}$, {[}CSN{]}$^{-}$ & igual & ambos \\
b & Trigonal plana: BF$_{3}$ (SO$_{3}$ também serve); pirâmide trigonal: NH$_{3}$ (NF$_{3}$ também serve) & Trigonal plana: BF$_{3}$, BH$_{3}$ ou SO$_{3}$; pirâmide trigonal: NH$_{3}$ ou NF$_{3}$ & igual & ambos \\
c & Nitração do fluorobenzeno (C$_{6}$H$_{5}$F) com HNO$_{3}$/H$_{2}$SO$_{4}$: 1{-}flúor{-}4{-}nitrobenzeno (para, principal) e 1{-}flúor{-}2{-}nitrobenzeno (orto, minoritário), mais H$_{2}$O; as duas estruturas estão desenhadas & Nitração do C$_{6}$H$_{5}$F: produtos principais o{-}fluoronitrobenzeno (1{-}fluoro{-}2{-}nitrobenzeno) e p{-}fluoronitrobenzeno (1{-}fluoro{-}4{-}nitrobenzeno) & equivalente & ambos \\
d & H$_{2}$S < CH$_{3}$OH < H$_{2}$O, ou seja, (ii) < (iii) < (i) & H$_{2}$S < CH$_{3}$OH < H$_{2}$O, ou seja, (ii) < (iii) < (i) & igual & ambos \\
\bottomrule\end{longtable}%
\begin{itemize}%
\item%
\textbf{Identificação dos átomos)} Identificação idêntica. A contagem de elétrons de cada configuração dá os números atômicos 1, 5, 6, 7, 8, 9 e 16.%
\item%
\textbf{a)} Os dois chegam a 3 arranjos lineares, e o critério é o mesmo: o átomo central define o arranjo, já que a cadeia lida de trás para a frente é a mesma. A resolução descarta o arranjo cíclico de 3 membros e, além do que se pede, discute cargas formais. Refiz as cargas formais da tabela e todas estão certas: S=C=N (0,0,-1), S--C$\equiv$N (-1,0,0), C=N=S (-2,+1,0), C=S=N (-2,+2,-1).%
\item%
\textbf{b)} O item pede um exemplo de cada geometria. A resolução dá BF$_{3}$ e NH$_{3}$, que estão na lista do gabarito, e cita como alternativas SO$_{3}$ e NF$_{3}$, que o gabarito também aceita. A justificativa por VSEPR está correta: 3 domínios sem par isolado no B e 4 domínios com 1 par isolado no N.%
\item%
\textbf{c)} Os dois apresentam e desenham os mesmos produtos, orto e para, com o F como orientador orto/para. A diferença é só de ênfase. O gabarito chama os dois de 'principais'. A resolução distingue o para como principal e o orto como minoritário, o que está de acordo com os dados experimentais: a nitração do fluorobenzeno dá cerca de 85--90\% de para e cerca de 10--13\% de orto. Conferi o desenho em TikZ: F em 90$^\circ$ e NO$_{2}$ em 270$^\circ$ dá para; F em 90$^\circ$ e NO$_{2}$ em 30$^\circ$ (vértice vizinho) dá orto. Está correto.%
\item%
\textbf{d)} Ordem idêntica, e a justificativa é a mesma: o H$_{2}$S só tem dipolo permanente e dispersão; H$_{2}$O e CH$_{3}$OH fazem ligações de hidrogênio, e a H$_{2}$O faz mais delas por molécula. Os pontos de ebulição citados ($\approx$ -60, 65 e 100 $^\circ$C) estão corretos.%
\end{itemize}%
\paragraph{Pontos de dúvida}%
\begin{itemize}%
\item%
\textit{(key)} Item c: o gabarito chama orto e para de 'produtos principais', mas experimentalmente o para predomina amplamente (cerca de 85--90\%). A resolução desenhou os dois e indicou o orto como minoritário. Vale confirmar com o professor se a banca exige os dois como 'principais' ou aceita a distinção. Na prática a nota não deve mudar, porque as duas estruturas foram desenhadas.%
\item%
\textit{(statement)} Item a: o enunciado não define 'arranjo atômico'. A resolução e o gabarito contam só cadeias lineares (3). Uma leitura que aceitasse um anel de 3 membros chegaria a 4, mas isso é pouco razoável quimicamente, e a resolução descarta o caso de forma explícita.%
\item%
\textit{(solution)} Item a: para o arranjo C--N--S, a tabela mostra só a estrutura {[}C=N=S{]}$^{-}$ (-2,+1,0). A alternativa {[}C$\equiv$N--S{]}$^{-}$ (-1,+1,-1), análoga à que costuma ser preferida para o fulminato, não foi considerada. A discussão vai além do que se pede e não afeta a resposta (3 arranjos).%
\item%
\textit{(solution)} Item d: a resolução usa M(CH$_{3}$OH) = 32,04 g/mol. Com a tabela da prova, 12,01 + 4$\times$1,01 + 16,00 = 32,05 g/mol. A diferença é de arredondamento e não muda nada.%
\end{itemize}%
\section*{Questão 6}%
\href{https://poliedroresolve.sistemapoliedro.com.br/vestibulares/ita/2026/ita-2-fase-dia-2-29-10-2025-2-fase-dia-2/dissertativa-6-quimica-2-fase-dia-2-ita-2026}{Resolução do gabarito}\par\smallskip%
\begin{longtable}{@{}p{0.06\textwidth} p{0.33\textwidth} p{0.33\textwidth} p{0.1\textwidth} p{0.08\textwidth}@{}}\toprule Item & Resolução & Gabarito & Veredito & Certo \\ \midrule
a & pKa(H3O+) = -1,70; pKb(CH3CC-) = -11,00; K1 = 10\^{}-26,70 $\approx$ 2,0$\times$10$^{-}$$^2$$^{7}$ < 1, então não favorece os produtos & pKa(H3O+) = -1,70; pKb(CH3CC-) = -11; Keq = 10\^{}(-1,70-25) < 1, então não favorece os produtos (equilíbrio deslocado para a esquerda) & igual & ambos \\
b & pKa(NH4+) = 9,25; K2 = 10\^{}-15,75 $\approx$ 1,8$\times$10$^{-}$$^{1}$$^{6}$ < 1, então não favorece os produtos & pKa(NH4+) = 14 - 4,75 = 9,25; Keq = 10\^{}(9,25-25) < 1, então não favorece os produtos & igual & ambos \\
c & pKa(NH3) = 36,00; K3 = 10\^{}11 > 1, então favorece os produtos & pKa(NH3) = 14 - (-22) = 36; Keq = 10\^{}(36-25) = 10\^{}11 > 1, então favorece os produtos & igual & ambos \\
\bottomrule\end{longtable}%
\begin{itemize}%
\item%
\textbf{a)} Os valores e a conclusão coincidem. Refiz a conta: 14 - 15,70 = -1,70; 14 - 25 = -11; 10\^{}-26,70 = 1,995$\times$10$^{-}$$^2$$^{7}$. Os dois textos observam que a convenção do pKb(H2O) = 15,70 é discutível. Tomando a atividade da água igual a 1, K1 = 10$^{-}$$^2$$^{5}$, e a conclusão não muda.%
\item%
\textbf{b)} Os valores e a conclusão são idênticos. A conta foi conferida: 10\^{}-15,75 = 1,78$\times$10$^{-}$$^{1}$$^{6}$, o que confere com a checagem Ka$\cdot$Kb/Kw.%
\item%
\textbf{c)} Os valores e a conclusão são idênticos. A resolução trata o Na+ como íon espectador, que é a hipótese correta. Conferi também por Ka$\cdot$Kb/Kw: 10$^{-}$$^2$$^{5}$$\cdot$10$^2$$^2$/10$^{-}$$^{1}$$^{4}$ = 10$^{1}$$^{1}$.%
\end{itemize}%
\paragraph{Pontos de dúvida}%
\begin{itemize}%
\item%
\textit{(statement)} O dado pKb(H2O) = 15,70 pressupõe a convenção que inclui {[}H2O{]} $\approx$ 55,5 mol/L na constante. Na convenção termodinâmica, com atividade da água igual a 1, valem pKa(H3O+) = 0 e pKb(H2O) = 14 (J. Chem. Educ. 2017). O Poliedro critica o dado, e a resolução comenta a mesma questão. A conclusão do item a não muda em nenhuma das duas convenções.%
\item%
\textit{(solution)} Na observação do item a, a resolução diz que o fator 10\^{}-1,70 vale cerca de 1/55,5. Na verdade 10\^{}1,70 $\approx$ 50,1. O valor exato com 55,5 mol/L seria 10\^{}1,744, o que corresponde a pKb $\approx$ 15,74. A imprecisão é pequena e não altera nenhuma resposta.%
\end{itemize}%
\section*{Questão 7}%
\href{https://poliedroresolve.sistemapoliedro.com.br/vestibulares/ita/2026/ita-2-fase-dia-2-29-10-2025-2-fase-dia-2/dissertativa-7-quimica-2-fase-dia-2-ita-2026}{Resolução do gabarito}\par\smallskip%
\begin{longtable}{@{}p{0.06\textwidth} p{0.33\textwidth} p{0.33\textwidth} p{0.1\textwidth} p{0.08\textwidth}@{}}\toprule Item & Resolução & Gabarito & Veredito & Certo \\ \midrule
a & Meio ácido: cátodo 4H? + 4e$^{-}$ $\to$ 2H$_{2}$; ânodo 2H$_{2}$O $\to$ O$_{2}$ + 4H? + 4e$^{-}$. Também dá a alternativa em meio básico: cátodo 4H$_{2}$O + 4e$^{-}$ $\to$ 2H$_{2}$ + 4OH$^{-}$; ânodo 4OH$^{-}$ $\to$ O$_{2}$ + 2H$_{2}$O + 4e$^{-}$. Global: 2H$_{2}$O $\to$ 2H$_{2}$ + O$_{2}$. & Cátodo: 2H$_{2}$O + 2e$^{-}$ $\to$ H$_{2}$ + 2OH$^{-}$; ânodo: 2H$_{2}$O $\to$ O$_{2}$ + 4H? + 4e$^{-}$ & equivalente & ambos \\
b & E = 2$\cdot$N\_A$\cdot$hc/$\lambda$ $\approx$ 5,32$\times$10$^{5}$ J = 532 kJ por mol de H$_{2}$ (E do fóton = 4,42$\times$10$^{-}$$^{1}$$^{9}$ J) & E = 2$\cdot$6$\cdot$10$^2$$^3$$\cdot$6,63$\cdot$10$^{-}$$^3$$^{4}$$\cdot$3$\cdot$10$^{8}$/(450$\cdot$10$^{-}$$^{9}$) $\approx$ 5,304$\times$10$^{5}$ J = 530,4 kJ & equivalente & ambos \\
c & E = 10$\cdot$0,2$\cdot$3600 = 7200 J $\to$ 1,63$\times$10$^2$$^2$ fótons = 2,71$\times$10$^{-}$$^2$ mol e$^{-}$ $\to$ n(H$_{2}$) $\approx$ 1,35$\times$10$^{-}$$^2$ mol & 7200 J / 5,304$\times$10$^{5}$ J/mol $\approx$ 1,36$\times$10$^{-}$$^2$ mol & equivalente & ambos \\
d & n = 5650/2,02 = 2,80$\times$10$^3$ mol; E\_min = 2,80$\times$10$^3$ $\cdot$ 532 kJ $\approx$ 1,49$\times$10$^{9}$ J & n = 5650/2 = 2825 mol; E = 530,4 kJ/mol $\cdot$ 2825 = 1 498 380 kJ $\approx$ 1,5 GJ & equivalente & ambos \\
e & 286 kJ/mol de H$_{2}$ ÷ 2 mol de fótons = 143 kJ/mol de fótons $\to$ $\lambda$\_max $\approx$ 837 nm & 286 kJ por mol de H$_{2}$ (2 mol de fótons): $\lambda$ = 2$\cdot$6$\cdot$10$^2$$^3$$\cdot$6,63$\cdot$10$^{-}$$^3$$^{4}$$\cdot$3$\cdot$10$^{8}$/(286$\cdot$10$^3$) $\approx$ 835 nm & equivalente & ambos \\
\bottomrule\end{longtable}%
\begin{itemize}%
\item%
\textbf{a)} Nos dois casos o cátodo reduz e gera H$_{2}$ e o ânodo oxida e gera O$_{2}$. A resolução escreve as semirreações com o mesmo meio (ácido ou básico). O gabarito põe o cátodo em meio básico e o ânodo em meio ácido. Mesmo assim, multiplicando o cátodo por 2 e somando, H? + OH$^{-}$ viram água e sobra a global 2H$_{2}$O $\to$ 2H$_{2}$ + O$_{2}$. As duas respostas estão corretas.%
\item%
\textbf{b)} Método idêntico: 2 mol de e$^{-}$, portanto 2 mol de fótons por mol de H$_{2}$. Recalculei: 2$\cdot$6,02$\times$10$^2$$^3$$\cdot$6,63$\times$10$^{-}$$^3$$^{4}$$\cdot$3,0$\times$10$^{8}$/450$\times$10$^{-}$$^{9}$ = 5,3217$\times$10$^{5}$ J. A diferença de 0,3\% vem só de o gabarito arredondar N\_A para 6$\times$10$^2$$^3$, quando a prova dá 6,02$\times$10$^2$$^3$. A resolução usa a constante fornecida e fica mais fiel ao enunciado.%
\item%
\textbf{c)} Recalculei: 7200/5,3217$\times$10$^{5}$ = 1,353$\times$10$^{-}$$^2$ mol com N\_A = 6,02$\times$10$^2$$^3$, e 7200/5,304$\times$10$^{5}$ = 1,357$\times$10$^{-}$$^2$ mol com N\_A = 6$\times$10$^2$$^3$. A diferença na terceira casa vem só do arredondamento de N\_A. Raciocínio igual nos dois.%
\item%
\textbf{d)} Recalculei. Com os dados da prova (M(H$_{2}$) = 2,02 g/mol e N\_A = 6,02$\times$10$^2$$^3$): n = 2797 mol e E = 2797 $\cdot$ 5,3217$\times$10$^{5}$ = 1,488$\times$10$^{9}$ J $\approx$ 1,49 GJ. Com M = 2 e N\_A = 6$\times$10$^2$$^3$: 2825 $\cdot$ 5,304$\times$10$^{5}$ = 1,498$\times$10$^{9}$ J $\approx$ 1,5 GJ. Os dois usam fótons de 450 nm (energia por fóton mínima que o semicondutor absorve). A diferença de 0,7\% vem só dos arredondamentos de constantes no gabarito.%
\item%
\textbf{e)} Mesmo raciocínio: a energia de 2 mol de fótons deve ser igual aos 286 kJ liberados por mol de H$_{2}$ na combustão. Recalculei: $\lambda$ = 2$\cdot$6,02$\times$10$^2$$^3$$\cdot$6,63$\times$10$^{-}$$^3$$^{4}$$\cdot$3,0$\times$10$^{8}$/2,86$\times$10$^{5}$ = 8,373$\times$10$^{-}$$^{7}$ m $\approx$ 837 nm. Com N\_A = 6$\times$10$^2$$^3$ dá 834,5 $\approx$ 835 nm. A massa de 5,65 kg se cancela, então o resultado é independente de M(H$_{2}$). A diferença de 2 nm vem só de N\_A.%
\end{itemize}%
\paragraph{Pontos de dúvida}%
\begin{itemize}%
\item%
\textit{(key)} O gabarito do Poliedro usa N\_A = 6$\times$10$^2$$^3$ mol$^{-}$$^{1}$ e M(H$_{2}$) = 2 g/mol, mas a prova fornece N\_A = 6,02$\times$10$^2$$^3$ e M(H) = 1,01 g/mol. Por isso os números mudam um pouco: b) 530,4 kJ contra 532 kJ; c) 1,36$\times$10$^{-}$$^2$ contra 1,35$\times$10$^{-}$$^2$ mol; d) 1,50 contra 1,49 GJ; e) 835 contra 837 nm. Os valores da resolução seguem as constantes do enunciado. As duas versões devem ser aceitas.%
\item%
\textit{(key)} No item a, o gabarito põe o cátodo em meio básico (2H$_{2}$O + 2e$^{-}$ $\to$ H$_{2}$ + 2OH$^{-}$) e o ânodo em meio ácido (2H$_{2}$O $\to$ O$_{2}$ + 4H? + 4e$^{-}$). A soma ainda dá a global correta, porque H? + OH$^{-}$ viram água, mas por convenção se escreve o par no mesmo meio. A resolução faz isso e dá as duas versões, ácida e básica.%
\item%
\textit{(statement)} No item d, 'energia mínima \ldots{} utilizando o semicondutor' foi lida pelos dois como a energia dos fótons de 450 nm, o limite de absorção do semicondutor (532 kJ por mol de H$_{2}$). Alguém poderia ler como o mínimo termodinâmico ($\Delta$H = 286 kJ/mol, ou $\Delta$G). Como o item e trata exatamente esse caso ideal, a leitura adotada parece a pretendida.%
\item%
\textit{(statement)} No item e, o enunciado pede que a energia luminosa seja igual à energia liberada na combustão, ou seja, |$\Delta$H| = 286 kJ/mol de H$_{2}$ (água líquida). Não usa $\Delta$G$^\circ$. Gabarito e resolução seguem essa leitura.%
\end{itemize}%
\section*{Questão 8}%
\href{https://poliedroresolve.sistemapoliedro.com.br/vestibulares/ita/2026/ita-2-fase-dia-2-29-10-2025-2-fase-dia-2/dissertativa-8-quimica-2-fase-dia-2-ita-2026}{Resolução do gabarito}\par\smallskip%
\begin{longtable}{@{}p{0.06\textwidth} p{0.33\textwidth} p{0.33\textwidth} p{0.1\textwidth} p{0.08\textwidth}@{}}\toprule Item & Resolução & Gabarito & Veredito & Certo \\ \midrule
a & CH4(g) + 2 O2(g) $\to$ CO2(g) + 2 H2O(g); $\Delta$Hc = -800 kJ/mol de CH4 (com $\Delta$Hf de H2O(g) = -241 kJ/mol). Para 12 MJ, $\Delta$H = -12 MJ com 15 mol de CH4. & CH4(g) + 2 O2(g) $\to$ CO2(g) + 2 H2O(g), $\Delta$H = -800 kJ/mol; para o requisito I: 15 CH4 + 30 O2 $\to$ 15 CO2 + 30 H2O(g), $\Delta$H = -12 000 kJ & igual & ambos \\
b & m(CH4) = 240,75 g ($\approx$ 241 g); m(O2) = 960 g & m(CH4) = 240 g; m(O2) = 960 g & equivalente & ambos \\
c & $\approx$ 0,644 g de CH4(g) (E útil = 40$\cdot$240,75 = 9630 J; Q da pré{-}câmara = 32,1 kJ) & 0,64 g (Q útil = 240$\cdot$40 = 9,6 kJ; Q de combustão = 32 kJ) & equivalente & ambos \\
d & Q = 0,25\%$\cdot$12 MJ = 30 kJ; m\_min = 30000/{[}(5/3)$\cdot$90{]} = 200 g de CH4(l) & Q = 30 kJ = m$\cdot$(5/3)$\cdot$90 $\Rightarrow$ m = 200 g (mínima porque se usou a temperatura máxima de 188 K) & igual & ambos \\
\bottomrule\end{longtable}%
\begin{itemize}%
\item%
\textbf{a)} Mesma equação e mesmo $\Delta$H. Conferi: $\Delta$Hf(H2O,g) = -286 + 45 = -241 kJ/mol e $\Delta$H = {[}-393 + 2(-241){]} - (-75) = -800 kJ/mol. As duas citam os 15 mol e os -12 MJ. A resolução não escreve explicitamente a equação multiplicada por 15, mas dá $\Delta$H = -12 MJ para 15 mol, que é a mesma informação.%
\item%
\textbf{b)} A diferença está só na massa molar usada. A resolução usa M(CH4) = 12,01 + 4$\cdot$1,01 = 16,05 g/mol, que é o valor da tabela da prova, e chega a 15$\cdot$16,05 = 240,75 g. O gabarito arredonda para M = 16 g/mol e chega a 240 g. A diferença é de 0,3\%, coisa de arredondamento. O O2 dá 30$\cdot$32,00 = 960 g nas duas.%
\item%
\textbf{c)} O raciocínio é o mesmo: energia útil = 40 J/g $\times$ massa bombeada, dividida pela eficiência de 0,30, e depois a massa vem de 800 kJ/mol. Conferi: 9630/0,3 = 32 100 J e 32,1/800$\cdot$16,05 = 0,644 g. Com M = 16, o gabarito obtém 32 kJ e 0,64 g. A diferença vem só da massa molar do item b. A resolução também testa a leitura em que a massa da pré{-}câmara está dentro dos 240,75 g e obtém 0,642 g, praticamente o mesmo valor.%
\item%
\textbf{d)} Resultado idêntico: 0,0025$\cdot$1,2$\cdot$10\^{}7 = 3,0$\cdot$10\^{}4 J e 3,0$\cdot$10\^{}4/150 = 200 g. As duas justificam o 'mínimo' pelo aquecimento até o limite de 188 K e desprezam a mudança de fase, como manda o enunciado.%
\end{itemize}%
\paragraph{Pontos de dúvida}%
\begin{itemize}%
\item%
\textit{(key)} O gabarito usa M(CH4) = 16 g/mol, mas a tabela da prova dá C = 12,01 e H = 1,01, o que leva a M = 16,05 g/mol. Por isso o gabarito chega a 240 g, 32 kJ e 0,64 g, enquanto a resolução chega a 240,75 g, 32,1 kJ e 0,644 g. É diferença de arredondamento e não muda a nota.%
\item%
\textit{(statement)} No item c, o enunciado não diz se a massa bombeada pela turbobomba é só o CH4 da câmara principal (240 g) ou se inclui o CH4 da pré{-}câmara. Também não diz se o calor da pré{-}câmara entra nos 12 MJ. O gabarito e a resolução supõem que os 240 g vão para a câmara principal e que a pré{-}câmara é adicional. A resolução mostra que a outra leitura dá 0,642 g, sem diferença relevante.%
\item%
\textit{(solution)} A resolução diz que a queima da pré{-}câmara 'consome mais 2,57 g de O2'. Isso supõe, sem declarar, que esse O2 se soma aos 960 g do item b. O item b não pede esse valor e ele não afeta nenhuma resposta, mas a afirmação depende de interpretação.%
\end{itemize}%
\section*{Questão 9}%
\href{https://poliedroresolve.sistemapoliedro.com.br/vestibulares/ita/2026/ita-2-fase-dia-2-29-10-2025-2-fase-dia-2/dissertativa-9-quimica-2-fase-dia-2-ita-2026}{Resolução do gabarito}\par\smallskip%
\begin{longtable}{@{}p{0.06\textwidth} p{0.33\textwidth} p{0.33\textwidth} p{0.1\textwidth} p{0.08\textwidth}@{}}\toprule Item & Resolução & Gabarito & Veredito & Certo \\ \midrule
a & Cu$^2$?, K?, H? (H$_{3}$O?), SO$_{4}$$^2$$^{-}$, HSO$_{4}$$^{-}$, Cr$_{2}$O$_{7}$$^2$$^{-}$; traços de OH$^{-}$, HCrO$_{4}$$^{-}$ e CrO$_{4}$$^2$$^{-}$ & Principais: H? (H$_{3}$O?), SO$_{4}$$^2$$^{-}$, Cu$^2$?, K?, Cr$_{2}$O$_{7}$$^2$$^{-}$; em quantidades muito pequenas: OH$^{-}$, HSO$_{4}$$^{-}$, CrO$_{4}$$^2$$^{-}$ & equivalente & ambos \\
b & Migração iônica: o Cu$^2$? (azul) sobe pelo braço do catodo (Pt$_{1}$) e o Cr$_{2}$O$_{7}$$^2$$^{-}$ (alaranjado) sobe pelo braço do anodo (Pt$_{2}$). Isso separa as cores que, juntas, davam o marrom; os demais íons são incolores. & Migração iônica: Pt$_{1}$ (polo ?) atrai os cátions K?, H$_{3}$O? (incolores) e Cu$^2$? (azul), e por isso fica azul perto de Pt$_{1}$. Pt$_{2}$ (polo ?) atrai os ânions SO$_{4}$$^2$$^{-}$ (incolor) e Cr$_{2}$O$_{7}$$^2$$^{-}$ (alaranjado), e por isso fica alaranjado perto de Pt$_{2}$. & igual & ambos \\
c & Pt$_{1}$ = catodo: 2H? + 2e$^{-}$ $\to$ H$_{2}$. Pt$_{2}$ = anodo: 2H$_{2}$O $\to$ O$_{2}$ + 4H? + 4e$^{-}$. A resolução observa ainda que o Cu$^2$? seria reduzido de preferência se chegasse à placa. & Pt$_{1}$ = catodo (?, redução): 2H?(aq) + 2e$^{-}$ $\to$ H$_{2}$(g). Pt$_{2}$ = anodo (?, oxidação): 2H$_{2}$O(?) $\to$ O$_{2}$(g) + 4H?(aq) + 4e$^{-}$. As placas estão em contato com o H$_{2}$SO$_{4}$. & igual & ambos \\
d & P$_{1}$ negativo e P$_{2}$ positivo. A corrente convencional vai P$_{2}$ $\to$ Pt$_{2}$ $\to$ solução $\to$ Pt$_{1}$ $\to$ P$_{1}$. Os elétrons vão de P$_{1}$ para Pt$_{1}$ e de Pt$_{2}$ para P$_{2}$. & P$_{1}$ negativo e P$_{2}$ positivo. Elétrons (sentido real) de P$_{1}$ para P$_{2}$. Corrente convencional de P$_{2}$ para P$_{1}$ pelo circuito externo. & igual & ambos \\
\bottomrule\end{longtable}%
\begin{itemize}%
\item%
\textbf{a)} As duas listas trazem as mesmas espécies. A resolução acrescenta o HCrO$_{4}$$^{-}$, que de fato existe no equilíbrio do Cr(VI) em meio ácido. A única diferença é a classificação do HSO$_{4}$$^{-}$: o gabarito o põe entre os traços, a resolução entre as espécies principais. Fiz a conta com Ka2 $\approx$ 1,0$\times$10$^{-}$$^2$: em H$_{2}$SO$_{4}$ de 0,01 a 1 mol/L, o HSO$_{4}$$^{-}$ é 59--99\% do enxofre dissolvido (92\% a 0,1 mol/L). Ele só é minoritário em soluções muito diluídas (16\% a 10$^{-}$$^3$ mol/L). A classificação da resolução é, portanto, a mais defensável. O item pede apenas a lista de espécies, então a diferença de classificação não muda a resposta.%
\item%
\textbf{b)} O mecanismo é o mesmo (migração de cátions e ânions coloridos em sentidos opostos sob o campo elétrico) e a atribuição das cores aos braços também. A resolução explica com mais cuidado por que cada faixa tem uma só cor: no braço de cada eletrodo, o íon de carga oposta é empurrado para baixo, e só a frente de mesma cor avança para dentro do H$_{2}$SO$_{4}$ incolor.%
\item%
\textbf{c)} Eletrodos e semirreações são idênticos. A observação condicional sobre o Cu$^2$? não contradiz o gabarito, já que a Figura (B) mostra incolor a região das placas. As justificativas da resolução estão corretas: E$^\circ$(K?/K) $\approx$ -2,9 V; E$^\circ$(O$_{2}$/H$_{2}$O) = 1,23 V < E$^\circ$(S$_{2}$O$_{8}$$^2$$^{-}$/SO$_{4}$$^2$$^{-}$) = 2,01 V; o Cr já está em +6 e não pode ser oxidado.%
\item%
\textbf{d)} A polaridade e o sentido da corrente são os mesmos. A resolução detalha o caminho: no circuito externo à fonte a corrente sai de P$_{2}$, passa por Pt$_{2}$, pela solução e por Pt$_{1}$ e volta a P$_{1}$. Isso é exatamente o que o gabarito chama de 'P$_{2}$ para P$_{1}$'. O sentido dos elétrons, de P$_{1}$ a P$_{2}$ passando pela célula, também é o mesmo.%
\end{itemize}%
\paragraph{Pontos de dúvida}%
\begin{itemize}%
\item%
\textit{(key)} O gabarito põe o HSO$_{4}$$^{-}$ entre as espécies 'em quantidades muito pequenas'. Com Ka2 $\approx$ 10$^{-}$$^2$, porém, o HSO$_{4}$$^{-}$ é a espécie de enxofre predominante em H$_{2}$SO$_{4}$ diluído típico: \textasciitilde{}92\% a 0,1 mol/L e \textasciitilde{}59\% a 0,01 mol/L. Como o item só pede a lista e ambos incluem o HSO$_{4}$$^{-}$, isso não afeta a nota, mas a classificação do gabarito é discutível.%
\item%
\textit{(statement)} O enunciado não informa a concentração do H$_{2}$SO$_{4}$ nem o tempo de eletrólise. A afirmação de que só H? se reduz no catodo, e não Cu$^2$?, depende de o cobre ainda não ter chegado à placa, como sugere a Figura (B) com a região das placas incolor. As duas resoluções adotam essa leitura.%
\item%
\textit{(solution)} A resolução lista HCrO$_{4}$$^{-}$ como traço, e o gabarito não o lista. Em meio ácido, o HCrO$_{4}$$^{-}$ é a forma monomérica relevante do Cr(VI), e chamá{-}lo de 'traço' depende da concentração. É um acréscimo correto e não deve ser penalizado.%
\end{itemize}%
\section*{Questão 10}%
\href{https://poliedroresolve.sistemapoliedro.com.br/vestibulares/ita/2026/ita-2-fase-dia-2-29-10-2025-2-fase-dia-2/dissertativa-10-quimica-2-fase-dia-2-ita-2026}{Resolução do gabarito}\par\smallskip%
\begin{longtable}{@{}p{0.06\textwidth} p{0.33\textwidth} p{0.33\textwidth} p{0.1\textwidth} p{0.08\textwidth}@{}}\toprule Item & Resolução & Gabarito & Veredito & Certo \\ \midrule
a & X = 226 & X = 226 & igual & ambos \\
b & Z = 86 & Z = 86 & igual & ambos \\
c & 3 partículas alfa & 3 partículas alfa & igual & ambos \\
\bottomrule\end{longtable}%
\begin{itemize}%
\item%
\textbf{a)} Q{-}232 (Z = 89) sofre decaimento $\beta$$^{-}$ e dá G{-}232 (Z = 90), com N = 142. E é isótono de G e tem Z = 84, então A\_E = 84 + 142 = 226. A é isóbaro de E, logo X = 226. A resolução e o gabarito seguem o mesmo caminho.%
\item%
\textbf{b)} Pela conservação de carga com 3 emissões $\alpha$: 92 = Z + 2$\cdot$3, logo Z = 86 (radônio). A resolução confere que A (Z = 86) e E (Z = 84) são elementos diferentes, como exige a definição de isóbaros.%
\item%
\textbf{c)} Pela conservação do número de massa: 238 = 226 + 4n, logo n = 3. A resolução e o gabarito chegam ao mesmo resultado.%
\end{itemize}%
\paragraph{Pontos de dúvida}%
\begin{itemize}%
\item%
\textit{(statement)} O enunciado escreve "massa atômica igual a 238" quando o que se usa na conta é o número de massa. Tanto a resolução quanto o gabarito leem assim, que é o padrão nesse tipo de questão.%
\item%
\textit{(statement)} Tanto a resolução quanto o gabarito supõem que J só emite partículas $\alpha$, sem nenhum $\beta$ no meio. O enunciado ("sofre decaimentos de partículas alfa") apoia essa leitura. Sem essa hipótese, o item c não teria resposta única.%
\end{itemize}%
\end{document}